\documentclass[11pt]{article}

\usepackage[margin=1in]{geometry}

\title{\textbf{EE 5239 Course Project Proposal}\\
Optimization for Multimodal Sensor Fusion}
\author{Timothy Hess}
\date{Fall 2026}

\begin{document}

\maketitle

\section*{Project Description}

The goal of this project is to investigate the use of optimization methods for
multimodal sensor fusion. Many engineering systems use multiple sensors to estimate
a quantity or state that cannot be measured accurately using a single sensor. The
performance of these systems depends on how information from the available sensors
is combined. This project will investigate how optimization can be used to determine
an effective method of combining measurements from multiple sensors.

A publicly available dataset containing simultaneous measurements from multiple
sensors will be selected for the project. Preference will be given to physiological
sensor data, such as measurements related to cardiovascular, respiratory, or activity
signals. The selected data will be processed and divided into appropriate training
and testing datasets.

A sensor fusion model will then be developed to estimate a target quantity using the
available sensor measurements. The parameters of this model will be determined using
optimization. Different optimization methods will be implemented and compared to
investigate their convergence behavior, computational requirements, and resulting
estimation performance. The specific model and optimization algorithms will be
selected after reviewing the characteristics of the chosen dataset and relevant
literature.

The project will also investigate the importance of the individual sensors. Models
using the complete sensor set will be compared with models using reduced sensor sets
to determine whether similar estimation performance can be achieved with fewer
measurements. If appropriate for the selected model, regularization or other
optimization-based approaches to sensor selection will also be investigated.

The project will be implemented primarily in Python using standard numerical and
data-analysis libraries. The expected computational requirements should be suitable
for a personal computer. Results will be evaluated using appropriate estimation
error metrics as well as measures of optimization performance, such as convergence
behavior and computational time.

The expected outcome of the project is a comparison of optimization approaches for
multimodal sensor fusion and an analysis of the tradeoff between estimation
performance, computational requirements, and the number of sensors used. Additional
model complexity or alternative optimization formulations may be investigated if
time permits.

\end{document}